\section{Parameter tables}\label{app:params}

Table~\ref{tab:machines} gives the per-machine constants of the engine, generated from the engine's own modules: the classification constants of
Equation~\eqref{eq:k1k2}, the design throw that enters the specific energy of Equation~\eqref{eq:t10} and the throw term of $K_2$, the capacity
constants of Equation~\eqref{eq:capacity}, and the ranges of the setting and the speed over which the Latin-hypercube sweeps of the learned models
draw each machine class (the sweeps also draw the throw, $x_{63}$, $n$ and $A\times b$ over per-class ranges recorded in the surrogate contract). The
circuit constants are typical values of their class; only the secondary cone is replaced, on the calibrated lane, by the fitted relations of
Section~\ref{sec:m-hp500}. Table~\ref{tab:constants} collects the remaining constants of the engine, the evaluation and CrusherBayes, each with the module
that holds it; every one of them is named in the methods where it acts. A reimplementation that changes any of them changes the numbers of the report,
which is why they are listed with their exact values rather than rounded.

\begin{table}[!t]
\centering
\caption{Per-machine constants of the engine, generated from \texttt{classification.py}, \texttt{capacity.py}, \texttt{machines.py} and \texttt{sweep.py}:
classification constants, design throw, reference capacity, optimum speed, reference gap, and the setting and speed ranges of the learned-model sweeps.}
\label{tab:machines}
\scriptsize
% generated by figures/make_tables.py; do not edit by hand
\begin{tabular}{@{}l rrrrr r rrr rr@{}}
\toprule
Machine & $a_1$ & $b_1$ & $a_2$ & $b_2$ & $K_3$ & $T_{ref}$ & $Q_{ref}$ & $N_{opt}$ & $g_{ref}$ & CSS range & Speed range \\
 & (mm) &  & (mm) &  &  & (mm) & (t/h) & (rpm) & (mm) & (mm) & (rpm) \\
\midrule
gyratory & $0.6$ & $0.95$ & $10.0$ & $2.10$ & $2.0$ & $30$ & $4500$ & $150$ & $175$ & $120$--$240$ & $100$--$200$ \\
jaw & $2.0$ & $0.95$ & $20.0$ & $2.30$ & $2.1$ & $38$ & $700$ & $300$ & $110$ & $50$--$160$ & $180$--$380$ \\
secondary cone & $0.2$ & $0.90$ & $6.0$ & $1.95$ & $2.3$ & $30$ & $520$ & $360$ & $38$ & $6$--$90$ & $180$--$560$ \\
tertiary cone & $0.1$ & $0.85$ & $3.0$ & $1.75$ & $2.5$ & $16$ & $240$ & $400$ & $16$ & $4$--$22$ & $220$--$600$ \\
short-head cone & $0.1$ & $0.82$ & $2.4$ & $1.65$ & $2.7$ & $16$ & $200$ & $480$ & $12$ & $4$--$16$ & $380$--$600$ \\
\bottomrule
\end{tabular}

\end{table}

\begin{table}[!t]
\centering
\caption{Further constants of the engine, of the leave-one-survey-out evaluation and of CrusherBayes, with their values as the code holds them and the
module of \texttt{data-pipeline/\allowbreak{}pipeline/\allowbreak{}} that defines each.}
\label{tab:constants}
\scriptsize
\begin{tabular}{@{}l l l@{}}
\toprule
Constant & Value & Module \\
\midrule
Top of the sieve grid, number of classes, ratio & $256$~mm, $28$, $\sqrt2$ & \texttt{model/\allowbreak{}crusher/\allowbreak{}engine.py}, \texttt{sieve.py} \\
Readout sizes & $1$, $4$, $8$, $16$, $32$~mm & \texttt{model/\allowbreak{}crusher/\allowbreak{}engine.py} \\
Ill-conditioning threshold of $\kappa$ & $10^{6}$ & \texttt{model/\allowbreak{}crusher/\allowbreak{}engine.py} \\
Reduction ratio flagged as negligible & below $1.05$ & \texttt{model/\allowbreak{}crusher/\allowbreak{}engine.py} \\
Austin slopes $\gamma$, $\beta$ & $0.35$, $4.2$ & \texttt{model/\allowbreak{}crusher/\allowbreak{}breakage.py} \\
$A$; energy coefficient; reference speed & $60\%$; $0.033$~kWh/t per mm; $350$~rpm & \texttt{model/\allowbreak{}crusher/\allowbreak{}breakage.py} \\
Cap of $t_{10}$; clamp of $\varphi$ & $0.44$; $[0.02, 0.98]$ & \texttt{model/\allowbreak{}crusher/\allowbreak{}breakage.py} \\
Throw fraction of $K_2$; floor of $K_2$ & $0.5$; $1.05\,K_1$ & \texttt{model/\allowbreak{}crusher/\allowbreak{}classification.py} \\
Nip angle base (cones, gyratory; jaw); friction & $22^\circ$; $19^\circ$; $\mu_f=0.35$ & \texttt{model/\allowbreak{}crusher/\allowbreak{}capacity.py} \\
HP500 throw and speed; feed exponent; blend $W_i$ & $34$~mm, $360$~rpm; $0.9$; $9.0$~kWh/t & \texttt{model/\allowbreak{}crusher/\allowbreak{}hp500.py} \\
HP500 support (setting, feed $F_{80}$) & $38$ to $55$~mm, $47$ to $140$~mm & \texttt{model/\allowbreak{}crusher/\allowbreak{}envelope.py} \\
Lane budgets (run time, trace size) & $1500$~ms, $256$~KiB & \texttt{core/\allowbreak{}gate.py} \\
Units slack & one full range on each side & \texttt{core/\allowbreak{}unit\_gate.py} \\
Residual clip; ridge penalty; band factor; seed & $0.20$; $1.0$; $1.2816$; $42$ & \texttt{hp500/\allowbreak{}loo.py} \\
Morrell $K_s$; $K_2$ (closed, open circuit) & $55$; $1.0$, $1.19$ & \texttt{model/\allowbreak{}crusher/\allowbreak{}morrell.py} \\
Discrepancy prior sd $\tau_Q$, $\tau_P$ & $140$~t/h, $36$~kW & \texttt{hp500/\allowbreak{}crusherbayes.py} \\
Noise prior medians; log sd; log bounds & $60$~t/h, $15$~kW; $1$; $[\ln1, \ln1000]$ & \texttt{hp500/\allowbreak{}crusherbayes.py} \\
Walkers; steps; burn-in; thinning; draws & $20$; $7000$; $1000$; $10$; $20$ & \texttt{hp500/\allowbreak{}crusherbayes.py} \\
Moves (DE, DE-snooker); convergence factor; margin & $0.8$, $0.2$; $50$; $3$ points & \texttt{hp500/\allowbreak{}crusherbayes.py} \\
Drift projection grid; bisection steps; slack & $64$; $60$; $4$ ulp & \texttt{model/\allowbreak{}crusher/} (drift functions) \\
Learned models: epochs; learning rate; weight decay & $500$ and $400$; $2\times10^{-3}$; $10^{-5}$ & \texttt{stages/\allowbreak{}train.py} \\
Learned models: points per class; seeds; noise; opset & $1400$ and $480$; $12345$, $67890$; $0.05$; $17$ & \texttt{stages/\allowbreak{}train.py} \\
\bottomrule
\end{tabular}
\end{table}

\section{The integrated discrepancy of CrusherBayes}\label{app:derivation}

The marginal likelihood of Equation~\eqref{eq:cbmarginal} and the conditional of Equation~\eqref{eq:cbcond} follow from the Gaussian structure of the
model. For one target and one fold, write the residual of the engine on the $m$ training surveys as $\mathbf r=\mathbf y-f(\boldsymbol\vartheta)$. By
Equation~\eqref{eq:cbmodel}, $\mathbf r=\mathrm Z\mathbf d+\boldsymbol\varepsilon$ with $\mathbf d\sim N(\mathbf 0,\tau^2\mathrm I_3)$ independent of
$\boldsymbol\varepsilon\sim N(\mathbf 0,\sigma^2\mathrm I_m)$, so $\mathbf r$ is Gaussian with mean zero and covariance
\begin{equation}
\mathrm{Cov}[\mathbf r] = \mathrm Z\,\mathrm{Cov}[\mathbf d]\,\mathrm Z^{\top} + \mathrm{Cov}[\boldsymbol\varepsilon] = \tau^2\mathrm Z\mathrm Z^{\top} + \sigma^2\mathrm I_m = \boldsymbol\Sigma ,
\label{eq:app-cov}
\end{equation}
which is Equation~\eqref{eq:cbmarginal}. The pair $(\mathbf d,\mathbf r)$ is jointly Gaussian with $\mathrm{Cov}[\mathbf d,\mathbf r]=\tau^2\mathrm Z^{\top}$,
so the conditional of $\mathbf d$ given $\mathbf r$ is Gaussian with
\begin{equation}
\mathbb E[\mathbf d\mid\mathbf r] = \tau^2\mathrm Z^{\top}\boldsymbol\Sigma^{-1}\mathbf r,\qquad
\mathrm{Cov}[\mathbf d\mid\mathbf r] = \tau^2\mathrm I_3 - \tau^4\,\mathrm Z^{\top}\boldsymbol\Sigma^{-1}\mathrm Z ,
\label{eq:app-cond}
\end{equation}
which is Equation~\eqref{eq:cbcond}. The log density of Equation~\eqref{eq:cbmarginal} is evaluated with the Cholesky factor $\mathrm L$ of
$\boldsymbol\Sigma$ as $-\tfrac12\|\mathrm L^{-1}\mathbf r\|^2-\sum_k\ln L_{kk}-\tfrac m2\ln2\pi$, and a covariance that is not positive definite returns
minus infinity, so the sampler rejects that point. Because $\mathrm Z$ has three columns and $m=9$ or $10$ rows, $\boldsymbol\Sigma$ is a $9\times9$ or
$10\times10$ matrix, and the integration replaces six sampled dimensions per fold by one factorisation per target and posterior evaluation, which keeps each evaluation at the cost of the engine runs. The code checks
the marginal against a Monte Carlo of the model with the coefficients drawn explicitly.

\section{Contract schemas}\label{app:schemas}

The tables below are generated from the schema files of \texttt{docs/\allowbreak{}contracts/\allowbreak{}schemas/\allowbreak{}} and from the OpenAPI document by \texttt{make\_tables.py}. Each
row is one field: its dotted path (one level of nesting is expanded; \texttt{[]} marks the items of an array), its type, its range or allowed values as the
schema states them, whether it is required at its level, and its description, shortened where it is long. A field described only by a reference to a
shared definition carries the type of that definition. Rules that a schema cannot state, such as a count that must equal the length of a list or a path that
must exist, are checked by the producer and consumer tests of each contract and are described in Section~\ref{sec:data}; the units of each field are in its
description or in the subsection of that section. The catalog row of every contract, with its producer, its consumers, its swap test and its missing-data
behaviour, is the authority when this appendix and the catalog differ.

\subsubsection*{params-record}
% generated by figures/make_tables.py; do not edit by hand
\begingroup\scriptsize
\begin{longtable}{@{}p{0.22\textwidth} p{0.11\textwidth} p{0.2\textwidth} p{0.05\textwidth} p{0.34\textwidth}@{}}
\toprule
Field & Type & Range or values & Req. & Meaning \\
\midrule
\endhead
\bottomrule
\endfoot
\texttt{case\_id} & string & pattern \textasciicircum{}[A-Z][A-Z0-9-]\{1,15\}\$ & yes & The case identifier, for example S01 or HP500-01. \\
\texttt{machine} & enum & cone-sec, cone-tert, cone-short-head, gyratory, jaw & yes & The crusher kind; the engine's per-machine constants follow it. \\
\texttt{css\_mm} & number & > 0; <= 600 & yes & Closed-side setting, the narrowest gap between the mantle and the concave (or the jaws), mm. \\
\texttt{throw\_mm} & number & >= 1; <= 100 & yes & Eccentric throw, the open-side setting minus the closed-side setting, mm. \\
\texttt{speed\_rpm} & number & >= 50; <= 1000 & yes & Eccentric (gyration) speed, revolutions per minute. \\
\texttt{feed\_x63\_mm} & number & > 0; <= 2000 & yes & Rosin-Rammler characteristic size of the feed (63.2 percent passing), mm. \\
\texttt{feed\_m} & number & >= 0.3; <= 4 & yes & Rosin-Rammler uniformity exponent of the feed, dimensionless. \\
\texttt{ore\_axb} & number & >= 10; <= 1000 & yes & JKMRC drop-weight ore competence A times b, dimensionless; lower is harder. \\
\texttt{ore\_wi} & number & >= 2; <= 40 & yes & Bond work index of the ore, kWh per tonne. \\
\end{longtable}
\endgroup

\subsubsection*{hp500-survey}
% generated by figures/make_tables.py; do not edit by hand
\begingroup\scriptsize
\begin{longtable}{@{}p{0.21\textwidth} p{0.10\textwidth} p{0.19\textwidth} p{0.05\textwidth} p{0.33\textwidth}@{}}
\caption{Fields of the \texttt{hp500-survey} contract, version $1.0.0$, generated from its schema file \texttt{hp500-survey.schema.json}: the dotted path of each field, its type, its range or allowed values as the schema states them, whether it is required at its level, and its meaning.}\label{tab:schema-hp500-survey}\\
\toprule
Field & Type & Range or values & Req. & Meaning \\
\midrule
\endfirsthead
\multicolumn{5}{@{}l}{\footnotesize Table~\thetable{} (continued)}\\
\toprule
Field & Type & Range or values & Req. & Meaning \\
\midrule
\endhead
\bottomrule
\endfoot
\texttt{survey} & integer & >= 1; <= 10 & yes & The survey number as published. \\
\texttt{f80\_mm} & number & > 0; <= 400 & yes & Measured feed 80 percent passing size, mm. \\
\texttt{css\_mm} & number & > 0; <= 200 & yes & Measured closed-side setting, mm. \\
\texttt{feed\_tph} & number & > 0; <= 5000 & yes & Measured solid feed rate, tonnes per hour. \\
\texttt{power\_kw} & number & > 0; <= 1000 & yes & Crusher power estimated from the measured current, kW. \\
\end{longtable}
\endgroup

\subsubsection*{case-index}
% generated by figures/make_tables.py; do not edit by hand
\begingroup\scriptsize
\begin{longtable}{@{}p{0.22\textwidth} p{0.11\textwidth} p{0.2\textwidth} p{0.05\textwidth} p{0.34\textwidth}@{}}
\toprule
Field & Type & Range or values & Req. & Meaning \\
\midrule
\endhead
\bottomrule
\endfoot
\texttt{schema} & const & chancadem.index/v1 & yes & -- \\
\texttt{engine\_version} & string & pattern \textasciicircum{}\textbackslash{}d+\textbackslash{}.\textbackslash{}d\{2\}\textbackslash{}.\textbackslash{}d\{3\}\$ & yes & The product version that baked the index; older than VERSION means the bake is stale. \\
\texttt{n\_cases} & integer & >= 1 & yes & -- \\
\texttt{default\_case} & string & -- & yes & -- \\
\texttt{cases} & array of object & 1 to n items & yes & -- \\
\texttt{cases[].case\_id} & string & -- & yes & -- \\
\texttt{cases[].title} & object & -- & yes & -- \\
\texttt{cases[].category} & object & -- & yes & -- \\
\texttt{cases[].real\_or\_synthetic} & enum & real, synthetic & yes & real: one of the measured HP500 surveys; synthetic: a circuit case of the engine. \\
\texttt{cases[].manifest\_path} & string & pattern \textasciicircum{}manifests/[\textasciicircum{}/\textbackslash{}\textbackslash{}]+\textbackslash{}.json\$ & yes & -- \\
\texttt{files} & array of string & items pattern \textasciicircum{}(?!/)(?!.*\textbackslash{}.\textbackslash{}.)[\textasciicircum{}\textbackslash{}\textbackslash{}]+\textbackslash{}.(json|onnx)\$ & yes & The further committed files the site serves (the model files, the HP500 validation documents, the model metrics), relative to data/derived; ... \\
\end{longtable}
\endgroup

\subsubsection*{case-manifest}
% generated by figures/make_tables.py; do not edit by hand
\begingroup\scriptsize
\begin{longtable}{@{}p{0.20\textwidth} p{0.09\textwidth} p{0.16\textwidth} p{0.04\textwidth} p{0.07\textwidth} p{0.30\textwidth}@{}}
\caption{Fields of the \texttt{case-manifest} contract, version $1.1.1$, generated from its schema file \texttt{case-manifest.schema.json}: the dotted path of each field down to its leaves (\texttt{[]} an element of a list, \texttt{<key>} any key of a map; a structure met again is listed once, where it first appears), its type, its range or allowed values as the schema states them, whether it is required at its level, its unit and its meaning.}\label{tab:schema-case-manifest}\\
\toprule
Field & Type & Range or values & Req. & Unit & Meaning \\
\midrule
\endfirsthead
\multicolumn{6}{@{}l}{\footnotesize Table~\thetable{} (continued)}\\
\toprule
Field & Type & Range or values & Req. & Unit & Meaning \\
\midrule
\endhead
\bottomrule
\endfoot
\texttt{schema} & const & c\allowbreak{}h\allowbreak{}a\allowbreak{}n\allowbreak{}c\allowbreak{}a\allowbreak{}d\allowbreak{}e\allowbreak{}m\allowbreak{}.\allowbreak{}m\allowbreak{}a\allowbreak{}n\allowbreak{}i\allowbreak{}f\allowbreak{}e\allowbreak{}s\allowbreak{}t\allowbreak{}/\allowbreak{}v\allowbreak{}1 & yes & -- & The document schema identifier, chancadem.manifest/v1. \\
\texttt{case\_\allowbreak{}id} & string & -- & yes & -- & The case identifier, for example S01 or HP500-01. \\
\texttt{title} & object & -- & yes & -- & The case title in English and Spanish. \\
\texttt{title.\allowbreak{}en} & string & -- & yes & -- & The English text. \\
\texttt{title.\allowbreak{}es} & string & -- & yes & -- & The Spanish text. \\
\texttt{category} & object & -- & yes & -- & The case category in English and Spanish. Its fields are those of \texttt{title}. \\
\texttt{stage} & enum & primary, secondary, tertiary, control & yes & -- & The circuit stage the case belongs to. \\
\texttt{control} & enum & null, negative, invalid, calibration & yes & -- & For a control case, what it controls; null otherwise. \\
\texttt{real\_\allowbreak{}or\_\allowbreak{}synthetic} & enum & real, synthetic & yes & -- & real: one of the measured HP500 surveys; synthetic: a circuit case of the engine. \\
\texttt{expected\_\allowbreak{}band} & object & -- & yes & -- & What a crushing engineer should see, in words; expect states the same in numbers. Its fields are those of \texttt{title}. \\
\texttt{context} & object & -- & yes & -- & The case write-up shown in the workbench Context: the question, the variables that matter and how to read the views. \\
\texttt{context.\allowbreak{}problem} & object & -- & yes & -- & The question the case answers. Its fields are those of \texttt{title}. \\
\texttt{context.\allowbreak{}variables} & object & -- & yes & -- & The variables that matter for it. Its fields are those of \texttt{title}. \\
\texttt{context.\allowbreak{}reading} & object & -- & yes & -- & How to read the views on it. Its fields are those of \texttt{title}. \\
\texttt{engine} & object & -- & yes & -- & The engine that baked the case. \\
\texttt{engine.\allowbreak{}package} & const & p\allowbreak{}i\allowbreak{}p\allowbreak{}e\allowbreak{}l\allowbreak{}i\allowbreak{}n\allowbreak{}e\allowbreak{}.\allowbreak{}m\allowbreak{}o\allowbreak{}d\allowbreak{}e\allowbreak{}l\allowbreak{}.\allowbreak{}c\allowbreak{}r\allowbreak{}u\allowbreak{}s\allowbreak{}h\allowbreak{}e\allowbreak{}r & yes & -- & The Python package of the canonical engine. \\
\texttt{engine.\allowbreak{}version} & string & pattern \textasciicircum{}\allowbreak{}\textbackslash{}\allowbreak{}d\allowbreak{}+\allowbreak{}\textbackslash{}\allowbreak{}.\allowbreak{}\textbackslash{}\allowbreak{}d\allowbreak{}\{\allowbreak{}2\allowbreak{}\}\allowbreak{}\textbackslash{}\allowbreak{}.\allowbreak{}\textbackslash{}\allowbreak{}d\allowbreak{}\{\allowbreak{}3\allowbreak{}\}\allowbreak{}\$ & yes & -- & The product version of the engine. \\
\texttt{engine.\allowbreak{}model} & string & -- & yes & -- & The models the engine chains, in words, with the HP500 calibration named on a calibrated case. \\
\texttt{params} & object & -- & yes & -- & The operating point the case was baked at (contract 1 accepted it); calibrated is true only on the HP500 surveys. \\
\texttt{params.\allowbreak{}machine} & enum & cone-sec, cone-tert, cone-short-head, gyratory, jaw & yes & -- & The crusher kind; the engine's per-machine constants follow it. \\
\texttt{params.\allowbreak{}css\_\allowbreak{}mm} & number & -- & yes & mm & Closed-side setting, the narrowest gap between the mantle and the concave (or the jaws) \\
\texttt{params.\allowbreak{}throw\_\allowbreak{}mm} & number & -- & yes & mm & Eccentric throw, the open-side setting minus the closed-side setting. \\
\texttt{params.\allowbreak{}speed\_\allowbreak{}rpm} & number & -- & yes & rpm & Eccentric (gyration) speed. \\
\texttt{params.\allowbreak{}feed\_\allowbreak{}x63\_\allowbreak{}mm} & number & -- & yes & mm & Rosin-Rammler characteristic size of the feed (63.2 percent passing) \\
\texttt{params.\allowbreak{}feed\_\allowbreak{}m} & number & -- & yes & -- & Rosin-Rammler uniformity exponent of the feed, dimensionless. \\
\texttt{params.\allowbreak{}ore\_\allowbreak{}axb} & number & -- & yes & -- & JKMRC drop-weight ore competence A times b, dimensionless; lower is harder. \\
\texttt{params.\allowbreak{}ore\_\allowbreak{}wi} & number & -- & yes & kWh/t & Bond work index of the ore. \\
\texttt{params.\allowbreak{}calibrated} & boolean & -- & yes & -- & True only on the HP500 surveys, whose engine runs with the calibration of Rocha et al. 2024. \\
\texttt{measured} & object & -- & no & -- & Only on a real case: the survey row exactly as contract hp500-survey accepted it. \\
\texttt{measured.\allowbreak{}survey} & integer & >= 1 & yes & -- & The survey number as published. \\
\texttt{measured.\allowbreak{}f80\_\allowbreak{}mm} & number & -- & yes & mm & Measured feed 80 percent passing size. \\
\texttt{measured.\allowbreak{}css\_\allowbreak{}mm} & number & -- & yes & mm & Measured closed-side setting. \\
\texttt{measured.\allowbreak{}feed\_\allowbreak{}tph} & number & -- & yes & t/h & Measured solid feed rate. \\
\texttt{measured.\allowbreak{}power\_\allowbreak{}kw} & number & -- & yes & kW & Crusher power estimated from the measured current. \\
\texttt{seed} & integer & -- & yes & -- & The seed the case was baked with; the manifest is a pure function of the operating point and this seed. \\
\texttt{artifact} & object & -- & yes & -- & The case's trace file. \\
\texttt{artifact.\allowbreak{}path} & string & pattern \textasciicircum{}\allowbreak{}(\allowbreak{}?\allowbreak{}!\allowbreak{}/\allowbreak{})\allowbreak{}(\allowbreak{}?\allowbreak{}!\allowbreak{}.\allowbreak{}*\allowbreak{}\textbackslash{}\allowbreak{}.\allowbreak{}\textbackslash{}\allowbreak{}.\allowbreak{})\allowbreak{}[\allowbreak{}\textasciicircum{}\allowbreak{}\textbackslash{}\allowbreak{}\textbackslash{}\allowbreak{}]\allowbreak{}+\allowbreak{}\textbackslash{}\allowbreak{}.\allowbreak{}j\allowbreak{}s\allowbreak{}o\allowbreak{}n\allowbreak{}\$ & yes & -- & Relative to data/derived; no leading slash, no '..', no backslash. \\
\texttt{artifact.\allowbreak{}format} & const & j\allowbreak{}s\allowbreak{}o\allowbreak{}n & yes & -- & The file format, json. \\
\texttt{artifact.\allowbreak{}trace\_\allowbreak{}schema} & const & c\allowbreak{}h\allowbreak{}a\allowbreak{}n\allowbreak{}c\allowbreak{}a\allowbreak{}d\allowbreak{}e\allowbreak{}m\allowbreak{}.\allowbreak{}t\allowbreak{}r\allowbreak{}a\allowbreak{}c\allowbreak{}e\allowbreak{}/\allowbreak{}v\allowbreak{}1 & yes & -- & The schema identifier of the trace, chancadem.trace/v1. \\
\texttt{artifact.\allowbreak{}bytes} & integer & >= 1 & yes & bytes & The size of the trace file as committed. \\
\texttt{artifact.\allowbreak{}sha256} & string & pattern \textasciicircum{}\allowbreak{}[\allowbreak{}0\allowbreak{}-\allowbreak{}9\allowbreak{}a\allowbreak{}-\allowbreak{}f\allowbreak{}]\allowbreak{}\{\allowbreak{}6\allowbreak{}4\allowbreak{}\}\allowbreak{}\$ & yes & -- & SHA-256 of the trace file bytes as committed. \\
\texttt{lane} & enum & live, precompute & yes & -- & How the App serves the case: live (the engine may run it in the browser) or precompute (the App replays the committed trace). \\
\texttt{gate} & object & -- & yes & -- & The measured lane decision of data-pipeline/pipeline/core/gate.py classify\_lane: the engine run is timed in the bake (the median of several warmed runs) and judged against run\_ms\_budget, but the measured time is not stored, so the manifest stays a pure function of (params, seed). \\
\texttt{gate.\allowbreak{}lane} & enum & live, precompute & yes & -- & The lane the measurement assigns, the same as lane. \\
\texttt{gate.\allowbreak{}pure\_\allowbreak{}python} & boolean & -- & yes & -- & Whether the engine is pure Python. \\
\texttt{gate.\allowbreak{}wheels} & array of string & -- & yes & -- & The compiled packages the engine depends on; a live case may use only numpy. \\
\texttt{gate.\allowbreak{}trace\_\allowbreak{}bytes} & integer & >= 1 & yes & bytes & The size of the trace. \\
\texttt{gate.\allowbreak{}run\_\allowbreak{}ms\_\allowbreak{}budget} & number & > 0 & yes & ms & The time one engine run may take for a live case, 1500 ms. \\
\texttt{gate.\allowbreak{}trace\_\allowbreak{}bytes\_\allowbreak{}budget} & integer & >= 1 & yes & bytes & The largest trace a live case may have, 262,144 bytes. \\
\texttt{gate.\allowbreak{}reasons} & array of string & -- & yes & -- & Why the case is not live; empty for a live case. \\
\texttt{flags} & array of enum & items css\_ge\_feed\_top, pass\_through, negligible\_reduction, ill\_conditioned, css\_nonpositive, feed\_nonpositive, f80\_lt\_css, css\_outside\_calibrated, f80\_outside\_calibrated, css\_below\_vendor\_minimum, reduction\_ratio\_above\_typical & yes & -- & Stable codes of the engine's flags and envelope codes for this case, in the engine's order: css\_ge\_feed\_top (nothing is gripped), pass\_through (CSS at or above F80), negligible\_reduction (reduction ratio near 1), ill\_conditioned (the solve's conditioning guard), the envelope codes css\_nonpositive, feed\_nonpositive, f80\_lt\_css, css\_outside\_calibrated, f80\_outside\_calibrated, and two sourced warnings that never make a case invalid: css\_below\_vendor\_minimum (the setting is below the vendor minimum of the machine type) and reduction\_ratio\_above\_typical (F80/P80 is above the typical ceiling of one stage). The web maps each code to its bilingual text. \\
\texttt{metrics} & object & -- & yes & -- & The case's result scalars (p80\_mm, p50\_mm, p20\_mm and f80\_mm in mm, reduction\_ratio, throughput\_tph in t/h, power\_kw in kW, specific\_energy\_kwh\_t in kWh/t) and, on a real case, the model error against the survey (throughput\_err\_pct and power\_err\_pct, in percent). \\
\texttt{models} & object & -- & yes & -- & Per model slot (catalog row id), what answers in this bake: the released ONNX model (real) or the slot's baseline. \\
\texttt{expect} & object & -- & yes & -- & The checked expectation of the case: [low, high] per result metric; the bake fails outside it. \\
\texttt{source} & object & -- & yes & -- & Where the case comes from: the authored circuit matrix, or the publication of a measured survey. \\
\texttt{source.\allowbreak{}kind} & const & a\allowbreak{}u\allowbreak{}t\allowbreak{}h\allowbreak{}o\allowbreak{}r\allowbreak{}e\allowbreak{}d or p\allowbreak{}u\allowbreak{}b\allowbreak{}l\allowbreak{}i\allowbreak{}c\allowbreak{}a\allowbreak{}t\allowbreak{}i\allowbreak{}o\allowbreak{}n & yes & -- & authored: a case of the circuit matrix. \\
\texttt{source.\allowbreak{}citation} & string & -- & no & -- & The publication the survey is transcribed from. \\
\texttt{source.\allowbreak{}doi} & string & pattern \textasciicircum{}\allowbreak{}1\allowbreak{}0\allowbreak{}\textbackslash{}\allowbreak{}.\allowbreak{}[\allowbreak{}0\allowbreak{}-\allowbreak{}9\allowbreak{}]\allowbreak{}\{\allowbreak{}4\allowbreak{},\allowbreak{}9\allowbreak{}\}\allowbreak{}/\allowbreak{}\textbackslash{}\allowbreak{}S\allowbreak{}+\allowbreak{}\$ & no & -- & Its DOI. \\
\texttt{source.\allowbreak{}license} & const & C\allowbreak{}C\allowbreak{} \allowbreak{}B\allowbreak{}Y\allowbreak{} \allowbreak{}4\allowbreak{}.\allowbreak{}0 & no & -- & The licence of the published data. \\
\end{longtable}
\endgroup

\subsubsection*{trace}
% generated by figures/make_tables.py; do not edit by hand
\begingroup\scriptsize
\begin{longtable}{@{}p{0.20\textwidth} p{0.09\textwidth} p{0.16\textwidth} p{0.04\textwidth} p{0.07\textwidth} p{0.30\textwidth}@{}}
\caption{Fields of the \texttt{trace} contract, version $1.0.2$, generated from its schema file \texttt{trace.schema.json}: the dotted path of each field down to its leaves (\texttt{[]} an element of a list, \texttt{<key>} any key of a map; a structure met again is listed once, where it first appears), its type, its range or allowed values as the schema states them, whether it is required at its level, its unit and its meaning.}\label{tab:schema-trace}\\
\toprule
Field & Type & Range or values & Req. & Unit & Meaning \\
\midrule
\endfirsthead
\multicolumn{6}{@{}l}{\footnotesize Table~\thetable{} (continued)}\\
\toprule
Field & Type & Range or values & Req. & Unit & Meaning \\
\midrule
\endhead
\bottomrule
\endfoot
\texttt{schema} & const & c\allowbreak{}h\allowbreak{}a\allowbreak{}n\allowbreak{}c\allowbreak{}a\allowbreak{}d\allowbreak{}e\allowbreak{}m\allowbreak{}.\allowbreak{}t\allowbreak{}r\allowbreak{}a\allowbreak{}c\allowbreak{}e\allowbreak{}/\allowbreak{}v\allowbreak{}1 & yes & -- & The document schema identifier, chancadem.trace/v1. \\
\texttt{case\_\allowbreak{}id} & string & -- & yes & -- & The case identifier, the same as its manifest's. \\
\texttt{sieve\_\allowbreak{}mm} & array of number & 2 to 200 items & yes & mm & Sieve aperture edges, coarse to fine (the engine grid: 256 mm, square-root-of-two series, 29 edges) \\
\texttt{mid\_\allowbreak{}mm} & array of number & 2 to 200 items & yes & mm & Geometric mid-size of each size class. \\
\texttt{feed\_\allowbreak{}passing} & array of number & 2 to 200 items & yes & fraction & Feed cumulative fraction passing at each edge. \\
\texttt{product\_\allowbreak{}passing} & array of number & 2 to 200 items & yes & fraction & Product cumulative fraction passing at each edge. \\
\texttt{feed\_\allowbreak{}mass} & array of number & 2 to 200 items & yes & fraction & Feed mass fraction in each class (sums to 1) \\
\texttt{product\_\allowbreak{}mass} & array of number & 2 to 200 items & yes & fraction & Product mass fraction in each class (sums to 1 within the mass-closure tolerance) \\
\texttt{classification} & array of number & 2 to 200 items & yes & fraction & Whiten classification function C(d) at each class mid-size: the probability a particle of that size is selected for breakage. \\
\texttt{appearance\_\allowbreak{}t10} & number & >= 0; <= 100 & yes & fraction & The t10 of the breakage appearance function used: the fraction of the progeny passing one tenth of the parent size, in [0, 0.44] (the engine caps it at 0.44) \\
\texttt{summary} & object & -- & yes & -- & The result scalars of the engine at the case's operating point. \\
\texttt{summary.\allowbreak{}f80\_\allowbreak{}mm} & number & >= 0 & yes & mm & Feed 80 percent passing size on the engine grid. \\
\texttt{summary.\allowbreak{}p80\_\allowbreak{}mm} & number & >= 0 & yes & mm & Product 80 percent passing size. \\
\texttt{summary.\allowbreak{}p50\_\allowbreak{}mm} & number & >= 0 & yes & mm & Product 50 percent passing size. \\
\texttt{summary.\allowbreak{}p20\_\allowbreak{}mm} & number & >= 0 & yes & mm & Product 20 percent passing size. \\
\texttt{summary.\allowbreak{}reduction\_\allowbreak{}ratio} & number & >= 0 & yes & -- & F80 over P80. \\
\texttt{summary.\allowbreak{}pct\_\allowbreak{}passing} & object & -- & yes & fraction & Product fraction passing at 1, 4, 8, 16 and 32 mm, keyed by the size. \\
\texttt{summary.\allowbreak{}throughput\_\allowbreak{}tph} & number & >= 0 & yes & t/h & Capacity of the machine at the operating point. \\
\texttt{summary.\allowbreak{}power\_\allowbreak{}kw} & number & >= 0 & yes & kW & Power draw. \\
\texttt{summary.\allowbreak{}specific\_\allowbreak{}energy\_\allowbreak{}kwh\_\allowbreak{}t} & number & >= 0 & yes & kWh/t & Specific comminution energy effectively applied. \\
\texttt{summary.\allowbreak{}oss\_\allowbreak{}mm} & number & >= 0 & yes & mm & Open-side setting, the closed-side setting plus the throw. \\
\texttt{summary.\allowbreak{}nip\_\allowbreak{}angle\_\allowbreak{}deg} & number & >= 0 & yes & deg & Nip angle between the mantle and the concave (or the jaws) at the setting. \\
\texttt{summary.\allowbreak{}nip\_\allowbreak{}limit\_\allowbreak{}deg} & number & >= 0 & yes & deg & Grip limit 2 arctan(mu) of the nip angle. \\
\texttt{summary.\allowbreak{}regime} & enum & choke, trickle, pass-through, invalid & yes & -- & The operating regime: choke, trickle, pass-through (setting at or above F80, or a reduction ratio near one) or invalid (setting at or above the feed's top size). \\
\texttt{summary.\allowbreak{}valid} & boolean & -- & yes & -- & False when the point is not valid: nothing gripped, an ill-conditioned breakage matrix or a physically implausible input. \\
\texttt{summary.\allowbreak{}out\_\allowbreak{}of\_\allowbreak{}envelope} & boolean & -- & yes & -- & True when any envelope check tripped (an implausible or an extrapolating input). \\
\texttt{summary.\allowbreak{}implausible} & boolean & -- & yes & -- & True for a physically nonsensical input, worse than an extrapolation. \\
\texttt{summary.\allowbreak{}mass\_\allowbreak{}closure} & number & >= 0 & yes & fraction & The absolute difference between the summed product and feed mass fractions, about 0. \\
\texttt{summary.\allowbreak{}cond\_\allowbreak{}estimate} & number or string & -- & yes & -- & Conditioning proxy of (I - B C); the string Infinity when the solve is singular. \\
\texttt{speed\_\allowbreak{}sweep} & object & -- & yes & -- & The capacity hump: the engine at the case's point with only the eccentric speed varied over the machine's range. \\
\texttt{speed\_\allowbreak{}sweep.\allowbreak{}speed\_\allowbreak{}rpm} & array of number & 2 to 200 items & yes & rpm & The eccentric speeds of the sweep. \\
\texttt{speed\_\allowbreak{}sweep.\allowbreak{}throughput\_\allowbreak{}tph} & array of number & 2 to 200 items & yes & t/h & The throughput at each speed; 0 where the engine cannot evaluate the point. \\
\texttt{speed\_\allowbreak{}sweep.\allowbreak{}power\_\allowbreak{}kw} & array of number & 2 to 200 items & yes & kW & The power at each speed; 0 where the engine cannot evaluate the point. \\
\texttt{css\_\allowbreak{}sweep} & object & -- & yes & -- & The setting trade-off: the engine at the case's point with only the closed-side setting varied. \\
\texttt{css\_\allowbreak{}sweep.\allowbreak{}css\_\allowbreak{}mm} & array of number & 2 to 200 items & yes & mm & The closed-side settings of the sweep. \\
\texttt{css\_\allowbreak{}sweep.\allowbreak{}p80\_\allowbreak{}mm} & array of number & 2 to 200 items & yes & mm & The product P80 at each setting; 0 where the engine cannot evaluate the point. \\
\texttt{css\_\allowbreak{}sweep.\allowbreak{}throughput\_\allowbreak{}tph} & array of number & 2 to 200 items & yes & t/h & The throughput at each setting; 0 where the engine cannot evaluate the point. \\
\texttt{css\_\allowbreak{}sweep.\allowbreak{}power\_\allowbreak{}kw} & array of number & 2 to 200 items & yes & kW & The power at each setting; 0 where the engine cannot evaluate the point. \\
\end{longtable}
\endgroup

\subsubsection*{hp500-backbone}
% generated by figures/make_tables.py; do not edit by hand
\begingroup\scriptsize
\begin{longtable}{@{}p{0.22\textwidth} p{0.11\textwidth} p{0.2\textwidth} p{0.05\textwidth} p{0.34\textwidth}@{}}
\toprule
Field & Type & Range or values & Req. & Meaning \\
\midrule
\endhead
\bottomrule
\endfoot
\texttt{schema} & const & chancadem.hp500.backbone/v1 & yes & -- \\
\texttt{source} & object & -- & yes & -- \\
\texttt{source.paper} & string & -- & yes & -- \\
\texttt{source.doi} & const & 10.3390/min14090919 & yes & -- \\
\texttt{source.license} & const & CC BY 4.0 & yes & -- \\
\texttt{source.table} & string & -- & yes & -- \\
\texttt{engine\_version} & string & pattern \textasciicircum{}\textbackslash{}d+\textbackslash{}.\textbackslash{}d\{2\}\textbackslash{}.\textbackslash{}d\{3\}\$ & yes & -- \\
\texttt{note} & string & -- & yes & -- \\
\texttt{n} & integer & >= 1 & yes & -- \\
\texttt{surveys} & array of object & 1 to n items & yes & -- \\
\texttt{surveys[].id} & string & pattern \textasciicircum{}HP500-[0-9]\{2\}\$ & yes & -- \\
\texttt{surveys[].n} & integer & >= 1 & yes & -- \\
\texttt{surveys[].measured} & object & -- & yes & -- \\
\texttt{surveys[].backbone} & object & -- & yes & -- \\
\end{longtable}
\endgroup

\subsubsection*{hp500-loo}
% generated by figures/make_tables.py; do not edit by hand
\begingroup\scriptsize
\begin{longtable}{@{}p{0.20\textwidth} p{0.09\textwidth} p{0.16\textwidth} p{0.04\textwidth} p{0.07\textwidth} p{0.30\textwidth}@{}}
\caption{Fields of the \texttt{hp500-loo} contract, version $1.2.1$, generated from its schema file \texttt{hp500-loo.schema.json}: the dotted path of each field down to its leaves (\texttt{[]} an element of a list, \texttt{<key>} any key of a map; a structure met again is listed once, where it first appears), its type, its range or allowed values as the schema states them, whether it is required at its level, its unit and its meaning.}\label{tab:schema-hp500-loo}\\
\toprule
Field & Type & Range or values & Req. & Unit & Meaning \\
\midrule
\endfirsthead
\multicolumn{6}{@{}l}{\footnotesize Table~\thetable{} (continued)}\\
\toprule
Field & Type & Range or values & Req. & Unit & Meaning \\
\midrule
\endhead
\bottomrule
\endfoot
\texttt{schema} & const & c\allowbreak{}h\allowbreak{}a\allowbreak{}n\allowbreak{}c\allowbreak{}a\allowbreak{}d\allowbreak{}e\allowbreak{}m\allowbreak{}.\allowbreak{}h\allowbreak{}p\allowbreak{}5\allowbreak{}0\allowbreak{}0\allowbreak{}.\allowbreak{}l\allowbreak{}o\allowbreak{}o\allowbreak{}/\allowbreak{}v\allowbreak{}1 & yes & -- & The document schema identifier, chancadem.hp500.loo/v1. \\
\texttt{source} & object & -- & yes & -- & The source publication of the surveys (Rocha, Campos, Silva and Tavares 2024), as the survey reader states it. \\
\texttt{source.\allowbreak{}paper} & string & -- & yes & -- & Authors, year and title of the source publication. \\
\texttt{source.\allowbreak{}doi} & const & 1\allowbreak{}0\allowbreak{}.\allowbreak{}3\allowbreak{}3\allowbreak{}9\allowbreak{}0\allowbreak{}/\allowbreak{}m\allowbreak{}i\allowbreak{}n\allowbreak{}1\allowbreak{}4\allowbreak{}0\allowbreak{}9\allowbreak{}0\allowbreak{}9\allowbreak{}1\allowbreak{}9 & yes & -- & The DOI of the source publication. \\
\texttt{source.\allowbreak{}license} & const & C\allowbreak{}C\allowbreak{} \allowbreak{}B\allowbreak{}Y\allowbreak{} \allowbreak{}4\allowbreak{}.\allowbreak{}0 & yes & -- & The licence of the source data, which the derived document carries. \\
\texttt{source.\allowbreak{}table} & string & -- & yes & -- & The table of the source publication the survey rows are transcribed from. \\
\texttt{engine\_\allowbreak{}version} & string & pattern \textasciicircum{}\allowbreak{}\textbackslash{}\allowbreak{}d\allowbreak{}+\allowbreak{}\textbackslash{}\allowbreak{}.\allowbreak{}\textbackslash{}\allowbreak{}d\allowbreak{}\{\allowbreak{}2\allowbreak{}\}\allowbreak{}\textbackslash{}\allowbreak{}.\allowbreak{}\textbackslash{}\allowbreak{}d\allowbreak{}\{\allowbreak{}3\allowbreak{}\}\allowbreak{}\$ & yes & -- & The product version that wrote the document (format X.XX.XXX); a version other than VERSION makes the row stale. \\
\texttt{n} & integer & >= 3 & yes & -- & The number of surveys, and of folds of each block. \\
\texttt{seed} & integer & -- & yes & -- & The seed of the label-shuffle permutation (seed for throughput, seed + 1 for power, seed + 2 for P80). \\
\texttt{config} & object & -- & yes & -- & The settings of the held-out evaluation. \\
\texttt{config.\allowbreak{}resid\_\allowbreak{}clip\_\allowbreak{}frac} & number & >= 0 & yes & fraction & The bound of the bounded residual: at most this fraction of the backbone value at the held-out survey. \\
\texttt{config.\allowbreak{}ridge\_\allowbreak{}lambda} & number & >= 0 & yes & -- & The ridge penalty on the two slopes of the bounded residual on the standardised features. \\
\texttt{config.\allowbreak{}features} & array of string & -- & yes & -- & The features of the residual and free-form fits, in order: the closed-side setting and f80. \\
\texttt{config.\allowbreak{}targets\_\allowbreak{}primary} & array of enum & items tph, kW & yes & -- & The measured targets: throughput (tph) and power (kW). \\
\texttt{config.\allowbreak{}target\_\allowbreak{}secondary} & array of const & items p\allowbreak{}8\allowbreak{}0 & yes & -- & The model output reported as self-consistency only: P80. \\
\texttt{config.\allowbreak{}morrell} & object & -- & no & -- & The constants and the fitted scalars of the two Morrell power variants. \\
\texttt{config.\allowbreak{}morrell.\allowbreak{}k\_\allowbreak{}s} & number & > 0 & yes & -- & K\_s of the hardness parameter S\_c (GMSG B5), 55 for a conventional crusher. \\
\texttt{config.\allowbreak{}morrell.\allowbreak{}k2} & number & > 0 & yes & -- & K2 of equation B4, 1.0 closed circuit and 1.19 open circuit. \\
\texttt{config.\allowbreak{}morrell.\allowbreak{}circuit} & enum & closed, open & yes & -- & The circuit the K2 constant is taken for. \\
\texttt{config.\allowbreak{}morrell.\allowbreak{}inputs} & string & -- & yes & -- & What the Morrell design column is computed from. \\
\texttt{config.\allowbreak{}morrell.\allowbreak{}fitted} & string & -- & yes & -- & Which scalars are fitted per fold and why. \\
\texttt{loo} & object & -- & yes & -- & The held-out evaluation per target. \\
\texttt{loo.\allowbreak{}tph} & object & -- & yes & t/h & The evaluation on the measured throughput. \\
\texttt{loo.\allowbreak{}tph.\allowbreak{}target} & enum & tph, kW, p80 & yes & -- & The target: tph, kW or p80. \\
\texttt{loo.\allowbreak{}tph.\allowbreak{}measured} & array of number & -- & yes & t/h, kW or mm & The values the predictions are graded against, in survey order: measured for tph and kW, the backbone's own P80 for p80. \\
\texttt{loo.\allowbreak{}tph.\allowbreak{}strict} & object & -- & yes & -- & Strict folds: every fitted scalar refit without the held-out survey; they also carry the predictive band (uq). \\
\texttt{loo.\allowbreak{}tph.\allowbreak{}strict.\allowbreak{}M0\_\allowbreak{}backbone} & object & -- & yes & -- & Scores of the calibrated backbone, with its scalars refit per fold (QT for throughput, zeta and Pn for power). \\ % internal-name-ok: schema field name
\texttt{loo.\allowbreak{}tph.\allowbreak{}strict.\allowbreak{}M0\_\allowbreak{}backbone.\allowbreak{}mape\_\allowbreak{}pct} & number & >= 0 & yes & \% & Mean absolute percentage error. \\ % internal-name-ok: schema field name
\texttt{loo.\allowbreak{}tph.\allowbreak{}strict.\allowbreak{}M0\_\allowbreak{}backbone.\allowbreak{}rmse} & number & >= 0 & yes & t/h, kW or mm & Root-mean-square error, in the target's unit. \\ % internal-name-ok: schema field name
\texttt{loo.\allowbreak{}tph.\allowbreak{}strict.\allowbreak{}M0\_\allowbreak{}backbone.\allowbreak{}r2} & number & -- & yes & -- & Coefficient of determination; 0 when the measured values have no variance. \\ % internal-name-ok: schema field name
\texttt{loo.\allowbreak{}tph.\allowbreak{}strict.\allowbreak{}M0\_\allowbreak{}backbone.\allowbreak{}maxdev\_\allowbreak{}pct} & number & >= 0 & yes & \% & The largest absolute percentage deviation over the surveys. \\ % internal-name-ok: schema field name
\texttt{loo.\allowbreak{}tph.\allowbreak{}strict.\allowbreak{}M1\_\allowbreak{}residual} & object & -- & yes & -- & Scores of the bounded residual: the backbone plus a ridge residual, clipped, on the standardised setting and f80. Its fields are those of \texttt{loo.\allowbreak{}tph.\allowbreak{}strict.\allowbreak{}M0\_\allowbreak{}backbone}. \\ % internal-name-ok: schema field name
\texttt{loo.\allowbreak{}tph.\allowbreak{}strict.\allowbreak{}M2\_\allowbreak{}freeform} & object & -- & yes & -- & Scores of the free-form fit, an unconstrained least-squares plane on the setting and f80 predicting the target. Its fields are those of \texttt{loo.\allowbreak{}tph.\allowbreak{}strict.\allowbreak{}M0\_\allowbreak{}backbone}. \\ % internal-name-ok: schema field name
\texttt{loo.\allowbreak{}tph.\allowbreak{}strict.\allowbreak{}constant\_\allowbreak{}mean} & object & -- & yes & -- & Scores of the control that predicts the mean of the fold's training surveys. Its fields are those of \texttt{loo.\allowbreak{}tph.\allowbreak{}strict.\allowbreak{}M0\_\allowbreak{}backbone}. \\
\texttt{loo.\allowbreak{}tph.\allowbreak{}strict.\allowbreak{}M1\_\allowbreak{}label\_\allowbreak{}shuffle} & object & -- & yes & -- & Scores of the control that fits the bounded residual to one seeded permutation of the measured values. Its fields are those of \texttt{loo.\allowbreak{}tph.\allowbreak{}strict.\allowbreak{}M0\_\allowbreak{}backbone}. \\ % internal-name-ok: schema field name
\texttt{loo.\allowbreak{}tph.\allowbreak{}strict.\allowbreak{}pred} & object & -- & yes & -- & The held-out prediction of each model at each survey, in survey order, rounded to two decimals. \\
\texttt{loo.\allowbreak{}tph.\allowbreak{}strict.\allowbreak{}pred.\allowbreak{}M0} & array of number & -- & yes & t/h, kW or mm & Predictions of the calibrated backbone. \\ % internal-name-ok: schema field name
\texttt{loo.\allowbreak{}tph.\allowbreak{}strict.\allowbreak{}pred.\allowbreak{}M1} & array of number & -- & yes & t/h, kW or mm & Predictions of the bounded residual. \\ % internal-name-ok: schema field name
\texttt{loo.\allowbreak{}tph.\allowbreak{}strict.\allowbreak{}pred.\allowbreak{}M2} & array of number & -- & yes & t/h, kW or mm & Predictions of the free-form fit. \\ % internal-name-ok: schema field name
\texttt{loo.\allowbreak{}tph.\allowbreak{}strict.\allowbreak{}pred.\allowbreak{}B\_\allowbreak{}lin} & array of number & -- & no & t/h or kW & Predictions of the setting-only least-squares line, throughput and power only. \\
\texttt{loo.\allowbreak{}tph.\allowbreak{}strict.\allowbreak{}pred.\allowbreak{}Morrell} & array of number & -- & no & kW & Predictions of Morrell through the origin, power only. \\
\texttt{loo.\allowbreak{}tph.\allowbreak{}strict.\allowbreak{}pred.\allowbreak{}Morrell\_\allowbreak{}noload} & array of number & -- & no & kW & Predictions of Morrell with a no-load term, power only. \\
\texttt{loo.\allowbreak{}tph.\allowbreak{}strict.\allowbreak{}uq} & object & -- & yes & -- & The 80 percent predictive band of the calibrated backbone from its strict held-out residuals, with its coverage. \\
\texttt{loo.\allowbreak{}tph.\allowbreak{}strict.\allowbreak{}uq.\allowbreak{}sigma} & number & >= 0 & yes & t/h, kW or mm & The sample standard deviation (ddof 1) of the ten held-out residuals of the calibrated backbone. \\
\texttt{loo.\allowbreak{}tph.\allowbreak{}strict.\allowbreak{}uq.\allowbreak{}band\_\allowbreak{}halfwidth} & number & >= 0 & yes & t/h, kW or mm & The half-width of the band, 1.2816 times sigma. \\
\texttt{loo.\allowbreak{}tph.\allowbreak{}strict.\allowbreak{}uq.\allowbreak{}nominal\_\allowbreak{}coverage} & number & >= 0; <= 1 & yes & fraction & The coverage the band is built for, 0.80. \\
\texttt{loo.\allowbreak{}tph.\allowbreak{}strict.\allowbreak{}uq.\allowbreak{}empirical\_\allowbreak{}coverage} & number & >= 0; <= 1 & yes & fraction & The share of surveys whose held-out residual lies within the half-width, in sample. \\
\texttt{loo.\allowbreak{}tph.\allowbreak{}strict.\allowbreak{}uq.\allowbreak{}band\_\allowbreak{}rel\_\allowbreak{}pct} & number & >= 0 & yes & \% & The half-width in percent of the mean measured value. \\
\texttt{loo.\allowbreak{}tph.\allowbreak{}strict.\allowbreak{}uq.\allowbreak{}note} & string & -- & yes & -- & How the band is built, in words. \\
\texttt{loo.\allowbreak{}tph.\allowbreak{}strict.\allowbreak{}uq.\allowbreak{}nested\_\allowbreak{}coverage} & number & >= 0; <= 1 & no & fraction & Out-of-sample coverage of the 80 percent band: the share of surveys whose held-out residual lies within the factor times the sigma of the other surveys' residuals. \\
\texttt{loo.\allowbreak{}tph.\allowbreak{}strict.\allowbreak{}uq.\allowbreak{}nested\_\allowbreak{}covered} & array of boolean & -- & no & -- & Per survey, in survey order: whether its residual lay within the nested band. \\
\texttt{loo.\allowbreak{}tph.\allowbreak{}strict.\allowbreak{}uq.\allowbreak{}nested\_\allowbreak{}sigma} & array of number & items >= 0 & no & t/h, kW or mm & Per survey, in survey order: the sample standard deviation of the held-out residuals of the other surveys. \\
\texttt{loo.\allowbreak{}tph.\allowbreak{}strict.\allowbreak{}uq.\allowbreak{}nested\_\allowbreak{}note} & string & -- & no & -- & States which coverage is in sample (empirical\_coverage) and which is out of sample (nested\_coverage). \\
\texttt{loo.\allowbreak{}tph.\allowbreak{}strict.\allowbreak{}B\_\allowbreak{}lin} & object & -- & no & -- & Scores of the setting-only least-squares line on the closed-side setting (throughput and power). Its fields are those of \texttt{loo.\allowbreak{}tph.\allowbreak{}strict.\allowbreak{}M0\_\allowbreak{}backbone}. \\
\texttt{loo.\allowbreak{}tph.\allowbreak{}strict.\allowbreak{}Morrell} & object & -- & no & -- & Scores of Morrell's size-specific energy through the origin, M\_ic refit per fold (power only). Its fields are those of \texttt{loo.\allowbreak{}tph.\allowbreak{}strict.\allowbreak{}M0\_\allowbreak{}backbone}. \\
\texttt{loo.\allowbreak{}tph.\allowbreak{}strict.\allowbreak{}Morrell\_\allowbreak{}noload} & object & -- & no & -- & Scores of Morrell's size-specific energy with a refit no-load power (power only). Its fields are those of \texttt{loo.\allowbreak{}tph.\allowbreak{}strict.\allowbreak{}M0\_\allowbreak{}backbone}. \\
\texttt{loo.\allowbreak{}tph.\allowbreak{}strict.\allowbreak{}morrell\_\allowbreak{}fit} & object & -- & no & -- & The scalars fitted in each fold (survey order, the held-out survey excluded from its own fold): M\_ic of the through-origin variant [kWh/t], M\_ic and the no-load power of the second variant, and whether that M\_ic is positive (a negative M\_ic is not Morrell's physics). \\
\texttt{loo.\allowbreak{}tph.\allowbreak{}strict.\allowbreak{}morrell\_\allowbreak{}fit.\allowbreak{}m\_\allowbreak{}ic\_\allowbreak{}kwh\_\allowbreak{}t} & array of number & -- & yes & kWh/t & M\_ic of the through-origin variant in each fold. \\
\texttt{loo.\allowbreak{}tph.\allowbreak{}strict.\allowbreak{}morrell\_\allowbreak{}fit.\allowbreak{}m\_\allowbreak{}ic\_\allowbreak{}noload\_\allowbreak{}kwh\_\allowbreak{}t} & array of number & -- & yes & kWh/t & M\_ic of the variant with a no-load term in each fold. \\
\texttt{loo.\allowbreak{}tph.\allowbreak{}strict.\allowbreak{}morrell\_\allowbreak{}fit.\allowbreak{}no\_\allowbreak{}load\_\allowbreak{}kw} & array of number & -- & yes & kW & The no-load power of that variant in each fold. \\
\texttt{loo.\allowbreak{}tph.\allowbreak{}strict.\allowbreak{}morrell\_\allowbreak{}fit.\allowbreak{}noload\_\allowbreak{}variant\_\allowbreak{}m\_\allowbreak{}ic\_\allowbreak{}positive} & array of boolean & -- & yes & -- & Whether that variant's M\_ic is positive in each fold (a negative M\_ic is not Morrell's physics). \\
\texttt{loo.\allowbreak{}tph.\allowbreak{}leaky} & object & -- & yes & -- & Leaky folds: every model fitted on all ten surveys, so the gap to the strict folds measures the optimism of a leaking evaluation. Its fields are those of \texttt{loo.\allowbreak{}tph.\allowbreak{}strict}. \\
\texttt{loo.\allowbreak{}kW} & object & -- & yes & kW & The evaluation on the measured power. Its fields are those of \texttt{loo.\allowbreak{}tph}. \\
\texttt{loo.\allowbreak{}p80} & object & -- & yes & mm & The self-consistency check on the backbone's own P80, a model output and not a measurement. Its fields are those of \texttt{loo.\allowbreak{}tph}. \\
\texttt{power\_\allowbreak{}parity} & object & -- & yes & -- & In-sample contrast of the power forms against the ten measured powers. \\
\texttt{power\_\allowbreak{}parity.\allowbreak{}note} & string & -- & yes & -- & What the contrast shows, in words. \\
\texttt{power\_\allowbreak{}parity.\allowbreak{}bond\_\allowbreak{}pd\_\allowbreak{}range\_\allowbreak{}kw} & array of number & exactly 2 items & yes & kW & The smallest and largest Bond draw of the engine over the surveys. \\
\texttt{power\_\allowbreak{}parity.\allowbreak{}measured\_\allowbreak{}range\_\allowbreak{}kw} & array of number & exactly 2 items & yes & kW & The smallest and largest measured power over the surveys. \\
\texttt{power\_\allowbreak{}parity.\allowbreak{}bare\_\allowbreak{}bond} & object & -- & yes & -- & The scores of the bare Bond draw against the measured power. Its fields are those of \texttt{loo.\allowbreak{}tph.\allowbreak{}strict.\allowbreak{}M0\_\allowbreak{}backbone}. \\
\texttt{power\_\allowbreak{}parity.\allowbreak{}engine\_\allowbreak{}pc\_\allowbreak{}1p30\_\allowbreak{}110} & object & -- & yes & -- & The scores of the paper's power form Pc = 1.30 Pd + 110 on the engine's Bond draw. Its fields are those of \texttt{loo.\allowbreak{}tph.\allowbreak{}strict.\allowbreak{}M0\_\allowbreak{}backbone}. \\
\texttt{power\_\allowbreak{}parity.\allowbreak{}refit\_\allowbreak{}pc} & object & -- & yes & -- & The scores of the same form refit by least squares on all ten surveys, with its two fitted scalars. \\
\texttt{power\_\allowbreak{}parity.\allowbreak{}refit\_\allowbreak{}pc.\allowbreak{}mape\_\allowbreak{}pct} & number & >= 0 & yes & \% & Mean absolute percentage error. \\
\texttt{power\_\allowbreak{}parity.\allowbreak{}refit\_\allowbreak{}pc.\allowbreak{}rmse} & number & >= 0 & yes & t/h, kW or mm & Root-mean-square error, in the target's unit. \\
\texttt{power\_\allowbreak{}parity.\allowbreak{}refit\_\allowbreak{}pc.\allowbreak{}r2} & number & -- & yes & -- & Coefficient of determination; 0 when the measured values have no variance. \\
\texttt{power\_\allowbreak{}parity.\allowbreak{}refit\_\allowbreak{}pc.\allowbreak{}maxdev\_\allowbreak{}pct} & number & >= 0 & yes & \% & The largest absolute percentage deviation over the surveys. \\
\texttt{power\_\allowbreak{}parity.\allowbreak{}refit\_\allowbreak{}pc.\allowbreak{}zeta} & number & -- & yes & -- & The fitted slope zeta of Pc = zeta Pd + Pn. \\
\texttt{power\_\allowbreak{}parity.\allowbreak{}refit\_\allowbreak{}pc.\allowbreak{}pn} & number & -- & yes & kW & The fitted no-load power Pn of Pc = zeta Pd + Pn. \\
\texttt{power\_\allowbreak{}parity.\allowbreak{}constant\_\allowbreak{}mean} & object & -- & yes & -- & The scores of the mean measured power used as the prediction for every survey. Its fields are those of \texttt{loo.\allowbreak{}tph.\allowbreak{}strict.\allowbreak{}M0\_\allowbreak{}backbone}. \\
\texttt{power\_\allowbreak{}parity.\allowbreak{}per\_\allowbreak{}survey} & array of object & -- & yes & -- & The four powers of each survey, in survey order. \\
\texttt{power\_\allowbreak{}parity.\allowbreak{}per\_\allowbreak{}survey[].\allowbreak{}n} & integer & -- & yes & -- & The survey number. \\
\texttt{power\_\allowbreak{}parity.\allowbreak{}per\_\allowbreak{}survey[].\allowbreak{}css} & number & -- & yes & mm & The measured closed-side setting. \\
\texttt{power\_\allowbreak{}parity.\allowbreak{}per\_\allowbreak{}survey[].\allowbreak{}measured} & number & -- & yes & kW & The measured power. \\
\texttt{power\_\allowbreak{}parity.\allowbreak{}per\_\allowbreak{}survey[].\allowbreak{}bare\_\allowbreak{}bond} & number & -- & yes & kW & The engine's Bond draw. \\
\texttt{power\_\allowbreak{}parity.\allowbreak{}per\_\allowbreak{}survey[].\allowbreak{}engine\_\allowbreak{}pc} & number & -- & yes & kW & The paper's form 1.30 Pd + 110. \\
\texttt{power\_\allowbreak{}parity.\allowbreak{}per\_\allowbreak{}survey[].\allowbreak{}refit\_\allowbreak{}pc} & number & -- & yes & kW & The form refit on all ten surveys. \\
\texttt{tph\_\allowbreak{}parity} & object & -- & yes & -- & In-sample contrast of the calibrated capacity backbone against the ten measured feed rates. \\
\texttt{tph\_\allowbreak{}parity.\allowbreak{}note} & string & -- & yes & -- & What the contrast shows, in words. \\
\texttt{tph\_\allowbreak{}parity.\allowbreak{}backbone} & object & -- & yes & -- & The scores of the backbone throughput against the measured feed rate. Its fields are those of \texttt{loo.\allowbreak{}tph.\allowbreak{}strict.\allowbreak{}M0\_\allowbreak{}backbone}. \\
\texttt{tph\_\allowbreak{}parity.\allowbreak{}per\_\allowbreak{}survey} & array of object & -- & yes & -- & The measured and the backbone throughput of each survey, in survey order. \\
\texttt{tph\_\allowbreak{}parity.\allowbreak{}per\_\allowbreak{}survey[].\allowbreak{}n} & integer & -- & yes & -- & The survey number. \\
\texttt{tph\_\allowbreak{}parity.\allowbreak{}per\_\allowbreak{}survey[].\allowbreak{}css} & number & -- & yes & mm & The measured closed-side setting. \\
\texttt{tph\_\allowbreak{}parity.\allowbreak{}per\_\allowbreak{}survey[].\allowbreak{}measured} & number & -- & yes & t/h & The measured feed rate. \\
\texttt{tph\_\allowbreak{}parity.\allowbreak{}per\_\allowbreak{}survey[].\allowbreak{}backbone} & number & -- & yes & t/h & The backbone throughput. \\
\end{longtable}
\endgroup

\subsubsection*{hp500-crusherbayes}
% generated by figures/make_tables.py; do not edit by hand
\begingroup\scriptsize
\begin{longtable}{@{}p{0.22\textwidth} p{0.11\textwidth} p{0.2\textwidth} p{0.05\textwidth} p{0.34\textwidth}@{}}
\toprule
Field & Type & Range or values & Req. & Meaning \\
\midrule
\endhead
\bottomrule
\endfoot
\texttt{schema} & const & chancadem.hp500.crusherbayes/v1 & yes & -- \\
\texttt{source} & object & -- & yes & -- \\
\texttt{source.paper} & string & -- & yes & -- \\
\texttt{source.doi} & const & 10.3390/min14090919 & yes & -- \\
\texttt{source.license} & const & CC BY 4.0 & yes & -- \\
\texttt{source.table} & string & -- & yes & -- \\
\texttt{engine\_version} & string & pattern \textasciicircum{}\textbackslash{}d+\textbackslash{}.\textbackslash{}d\{2\}\textbackslash{}.\textbackslash{}d\{3\}\$ & yes & -- \\
\texttt{n} & integer & >= 3 & yes & -- \\
\texttt{seed} & integer & -- & yes & -- \\
\texttt{inputs\_sha256} & string & pattern \textasciicircum{}[0-9a-f]\{64\}\$ & yes & SHA-256 over the survey table, the sampler size, the seed and the priors the document was sampled from. \\
\texttt{method} & object & -- & yes & -- \\
\texttt{method.sampler} & const & emcee & yes & -- \\
\texttt{method.sampler\_version} & string & -- & yes & -- \\
\texttt{method.move} & string & -- & yes & -- \\
\texttt{method.walkers} & integer & >= 2 & yes & -- \\
\texttt{method.steps} & integer & >= 1 & yes & -- \\
\texttt{method.burn\_in\_steps} & integer & >= 0 & yes & -- \\
\texttt{method.thin} & integer & >= 1 & yes & -- \\
\texttt{method.predictive\_draws\_per\_sample} & integer & >= 1 & yes & -- \\
\texttt{method.likelihood} & string & -- & yes & -- \\
\texttt{method.discrepancy\_prior\_sd} & object & -- & yes & -- \\
\texttt{method.noise\_prior} & object & -- & yes & -- \\
\texttt{method.p80} & string & -- & yes & -- \\
\texttt{method.interval\_80} & string & -- & yes & -- \\
\texttt{priors} & array of object & 1 to n items & yes & -- \\
\texttt{priors[].name} & string & -- & yes & -- \\
\texttt{priors[].unit} & string & -- & yes & -- \\
\texttt{priors[].nominal} & number & -- & yes & -- \\
\texttt{priors[].log} & boolean & -- & yes & -- \\
\texttt{priors[].mu} & number & -- & yes & -- \\
\texttt{priors[].sd} & number & > 0 & yes & -- \\
\texttt{priors[].lo} & number & -- & yes & -- \\
\texttt{priors[].hi} & number & -- & yes & -- \\
\texttt{priors[].source} & string & -- & yes & -- \\
\texttt{convergence} & object & -- & yes & -- \\
\texttt{convergence.criterion} & string & -- & yes & -- \\
\texttt{convergence.tau\_factor} & integer & >= 1 & yes & -- \\
\texttt{convergence.folds} & array of object & 2 to n items & yes & -- \\
\texttt{convergence.all\_converged} & boolean & -- & yes & -- \\
\texttt{loo} & object & -- & yes & -- \\
\texttt{loo.tph} & object & -- & yes & -- \\
\texttt{loo.kW} & object & -- & yes & -- \\
\texttt{comparators} & object & -- & yes & -- \\
\texttt{comparators.tph} & object & -- & yes & -- \\
\texttt{comparators.kW} & object & -- & yes & -- \\
\texttt{comparators\_note} & string & -- & yes & -- \\
\texttt{contraction} & array of object & 1 to n items & yes & -- \\
\texttt{contraction[].name} & string & -- & yes & -- \\
\texttt{contraction[].unit} & string & -- & yes & -- \\
\texttt{contraction[].nominal} & number & -- & yes & -- \\
\texttt{contraction[].prior\_sd} & number & > 0 & yes & -- \\
\texttt{contraction[].posterior\_mean} & number & -- & yes & -- \\
\texttt{contraction[].posterior\_sd} & number & >= 0 & yes & -- \\
\texttt{contraction[].contraction} & number & >= 0 & yes & -- \\
\texttt{contraction\_note} & string & -- & yes & -- \\
\texttt{kill\_criterion} & object & -- & yes & -- \\
\texttt{kill\_criterion.rule} & string & -- & yes & -- \\
\texttt{kill\_criterion.margin\_mape\_pts} & number & >= 0 & yes & -- \\
\texttt{kill\_criterion.coverage\_band} & array of number & 2 to 2 items & yes & -- \\
\texttt{kill\_criterion.power\_mape\_pct} & number & >= 0 & yes & -- \\
\texttt{kill\_criterion.power\_mape\_blin\_pct} & number & >= 0 & yes & -- \\
\texttt{kill\_criterion.power\_gain\_over\_blin\_pts} & number & -- & yes & -- \\
\texttt{kill\_criterion.a\_power\_beats\_blin\_by\_margin} & boolean & -- & yes & -- \\
\texttt{kill\_criterion.tph\_mape\_pct} & number & >= 0 & yes & -- \\
\texttt{kill\_criterion.tph\_mape\_m0\_pct} & number & >= 0 & yes & -- \\
\texttt{kill\_criterion.tph\_gain\_over\_m0\_pts} & number & -- & yes & -- \\
\texttt{kill\_criterion.a\_tph\_beats\_m0\_by\_margin} & boolean & -- & yes & -- \\
\texttt{kill\_criterion.a\_holds} & boolean & -- & yes & -- \\
\texttt{kill\_criterion.coverage} & object & -- & yes & -- \\
\texttt{kill\_criterion.coverage\_outside\_band} & object & -- & yes & -- \\
\texttt{kill\_criterion.mean\_width} & object & -- & yes & -- \\
\texttt{kill\_criterion.mean\_width\_crushercal} & object & -- & yes & -- \\
\texttt{kill\_criterion.band\_wider\_than\_crushercal} & object & -- & yes & -- \\
\texttt{kill\_criterion.b\_holds} & boolean & -- & yes & -- \\
\texttt{kill\_criterion.verdict} & enum & kept, dropped as a predictor & yes & -- \\
\end{longtable}
\endgroup

\subsubsection*{surrogate-model}
% generated by figures/make_tables.py; do not edit by hand
\begingroup\scriptsize
\begin{longtable}{@{}p{0.20\textwidth} p{0.09\textwidth} p{0.16\textwidth} p{0.04\textwidth} p{0.07\textwidth} p{0.30\textwidth}@{}}
\caption{Fields of the \texttt{surrogate-model} contract, version $1.0.1$, generated from its schema file \texttt{surrogate-model.schema.json}: the dotted path of each field down to its leaves (\texttt{[]} an element of a list, \texttt{<key>} any key of a map; a structure met again is listed once, where it first appears), its type, its range or allowed values as the schema states them, whether it is required at its level, its unit and its meaning.}\label{tab:schema-surrogate-model}\\
\toprule
Field & Type & Range or values & Req. & Unit & Meaning \\
\midrule
\endfirsthead
\multicolumn{6}{@{}l}{\footnotesize Table~\thetable{} (continued)}\\
\toprule
Field & Type & Range or values & Req. & Unit & Meaning \\
\midrule
\endhead
\bottomrule
\endfoot
\texttt{schema} & const & c\allowbreak{}h\allowbreak{}a\allowbreak{}n\allowbreak{}c\allowbreak{}a\allowbreak{}d\allowbreak{}e\allowbreak{}m\allowbreak{}.\allowbreak{}s\allowbreak{}u\allowbreak{}r\allowbreak{}r\allowbreak{}o\allowbreak{}g\allowbreak{}a\allowbreak{}t\allowbreak{}e\allowbreak{}/\allowbreak{}v\allowbreak{}1 & yes & -- & The contract schema identifier, chancadem.surrogate/v1. \\
\texttt{engine\_\allowbreak{}version} & string & pattern \textasciicircum{}\allowbreak{}\textbackslash{}\allowbreak{}d\allowbreak{}+\allowbreak{}\textbackslash{}\allowbreak{}.\allowbreak{}\textbackslash{}\allowbreak{}d\allowbreak{}\{\allowbreak{}2\allowbreak{}\}\allowbreak{}\textbackslash{}\allowbreak{}.\allowbreak{}\textbackslash{}\allowbreak{}d\allowbreak{}\{\allowbreak{}3\allowbreak{}\}\allowbreak{}\$ & yes & -- & The product version whose engine produced the training labels. \\
\texttt{machine\_\allowbreak{}order} & array of const & exactly 5 items & yes & -- & The machines in the order of the one-hot that opens the input vector. \\
\texttt{input\_\allowbreak{}order} & array of const & exactly 6 items & yes & -- & The continuous inputs in the order they follow the one-hot. \\
\texttt{in\_\allowbreak{}mean} & array of number & exactly 6 items & yes & -- & The mean of each continuous input on the training draw, in input\_order (the input is standardised as (value - in\_mean) / in\_std). \\
\texttt{in\_\allowbreak{}std} & array of number & exactly 6 items & yes & -- & The standard deviation of each continuous input on the training draw, in input\_order. \\
\texttt{output\_\allowbreak{}order} & array of const & exactly 10 items & yes & -- & The outputs in the order of the graph's output vector. \\
\texttt{out\_\allowbreak{}log} & array of boolean & exactly 10 items & yes & -- & Per output: true when the network predicts its log10 (P80, P50 and P20). \\
\texttt{out\_\allowbreak{}mean} & array of number & exactly 10 items & yes & -- & The mean of each output on the training draw (of its log10 where out\_log is true), in output\_order (the output is decoded as y * out\_std + out\_mean). \\
\texttt{out\_\allowbreak{}std} & array of number & exactly 10 items & yes & -- & The standard deviation of each output on the training draw (of the log10 where out\_log is true), in output\_order. \\
\texttt{graphs} & object & -- & yes & -- & The two exported graphs of the network. \\
\texttt{graphs.\allowbreak{}float32} & object & -- & yes & -- & The float32 graph, surrogate.onnx, the one the App runs. \\
\texttt{graphs.\allowbreak{}float32.\allowbreak{}path} & string & pattern \textasciicircum{}\allowbreak{}m\allowbreak{}o\allowbreak{}d\allowbreak{}e\allowbreak{}l\allowbreak{}s\allowbreak{}/\allowbreak{}[\allowbreak{}a\allowbreak{}-\allowbreak{}z\allowbreak{}0\allowbreak{}-\allowbreak{}9\allowbreak{}.\allowbreak{}-\allowbreak{}]\allowbreak{}+\allowbreak{}\textbackslash{}\allowbreak{}.\allowbreak{}o\allowbreak{}n\allowbreak{}n\allowbreak{}x\allowbreak{}\$ & yes & -- & Relative to data/derived. \\
\texttt{graphs.\allowbreak{}float32.\allowbreak{}bytes} & integer & >= 1; <= 5242880 & yes & bytes & The size of the graph file. \\
\texttt{graphs.\allowbreak{}float32.\allowbreak{}sha256} & string & pattern \textasciicircum{}\allowbreak{}[\allowbreak{}0\allowbreak{}-\allowbreak{}9\allowbreak{}a\allowbreak{}-\allowbreak{}f\allowbreak{}]\allowbreak{}\{\allowbreak{}6\allowbreak{}4\allowbreak{}\}\allowbreak{}\$ & yes & -- & SHA-256 of the graph file bytes. \\
\texttt{graphs.\allowbreak{}float32.\allowbreak{}opset} & integer & >= 13 & yes & -- & The ONNX operator set the graph is exported with. \\
\texttt{graphs.\allowbreak{}float32.\allowbreak{}input} & const & x & yes & -- & The name of the graph's input tensor, x. \\
\texttt{graphs.\allowbreak{}float32.\allowbreak{}output} & const & y & yes & -- & The name of the graph's output tensor, y. \\
\texttt{graphs.\allowbreak{}int8} & object & -- & yes & -- & The dynamically quantised graph, surrogate.int8.onnx. Its fields are those of \texttt{graphs.\allowbreak{}float32}. \\
\texttt{architecture} & string & -- & yes & -- & The layer sizes and activations, for example 11-64-64-32-10, SiLU. \\
\texttt{training} & object & -- & yes & -- & How the network was trained. \\
\texttt{training.\allowbreak{}seed} & integer & -- & yes & -- & The seed of the training draw and of torch. \\
\texttt{training.\allowbreak{}n\_\allowbreak{}train} & integer & >= 1 & yes & -- & The number of operating points the network was trained on. \\
\texttt{training.\allowbreak{}n\_\allowbreak{}test} & integer & >= 1 & yes & -- & The number of points of the independent held-out draw. \\
\texttt{training.\allowbreak{}sampler} & string & -- & yes & -- & How the operating points were drawn, with the bounds per input; the test draw is independent of the training draw. \\
\texttt{training.\allowbreak{}input\_\allowbreak{}ranges} & object & -- & yes & -- & The model's domain: the training range of every continuous input per machine. \\
\texttt{training.\allowbreak{}input\_\allowbreak{}ranges.\allowbreak{}cone-sec} & object & -- & yes & -- & The training ranges of the secondary cone. \\
\texttt{training.\allowbreak{}input\_\allowbreak{}ranges.\allowbreak{}cone-sec.\allowbreak{}css\_\allowbreak{}mm} & array of number & exactly 2 items & yes & mm & Closed-side setting range. \\
\texttt{training.\allowbreak{}input\_\allowbreak{}ranges.\allowbreak{}cone-sec.\allowbreak{}throw\_\allowbreak{}mm} & array of number & exactly 2 items & yes & mm & Eccentric throw range. \\
\texttt{training.\allowbreak{}input\_\allowbreak{}ranges.\allowbreak{}cone-sec.\allowbreak{}speed\_\allowbreak{}rpm} & array of number & exactly 2 items & yes & rpm & Eccentric speed range. \\
\texttt{training.\allowbreak{}input\_\allowbreak{}ranges.\allowbreak{}cone-sec.\allowbreak{}feed\_\allowbreak{}x63\_\allowbreak{}mm} & array of number & exactly 2 items & yes & mm & Rosin-Rammler characteristic feed size range. \\
\texttt{training.\allowbreak{}input\_\allowbreak{}ranges.\allowbreak{}cone-sec.\allowbreak{}feed\_\allowbreak{}m} & array of number & exactly 2 items & yes & -- & Rosin-Rammler uniformity exponent range. \\
\texttt{training.\allowbreak{}input\_\allowbreak{}ranges.\allowbreak{}cone-sec.\allowbreak{}ore\_\allowbreak{}axb} & array of number & exactly 2 items & yes & -- & Ore competence A times b range. \\
\texttt{training.\allowbreak{}input\_\allowbreak{}ranges.\allowbreak{}cone-tert} & object & -- & yes & -- & The training ranges of the tertiary cone. Its fields are those of \texttt{training.\allowbreak{}input\_\allowbreak{}ranges.\allowbreak{}cone-sec}. \\
\texttt{training.\allowbreak{}input\_\allowbreak{}ranges.\allowbreak{}jaw} & object & -- & yes & -- & The training ranges of the jaw crusher. Its fields are those of \texttt{training.\allowbreak{}input\_\allowbreak{}ranges.\allowbreak{}cone-sec}. \\
\texttt{training.\allowbreak{}input\_\allowbreak{}ranges.\allowbreak{}cone-short-head} & object & -- & yes & -- & The training ranges of the short-head cone. Its fields are those of \texttt{training.\allowbreak{}input\_\allowbreak{}ranges.\allowbreak{}cone-sec}. \\
\texttt{training.\allowbreak{}input\_\allowbreak{}ranges.\allowbreak{}gyratory} & object & -- & yes & -- & The training ranges of the gyratory crusher. Its fields are those of \texttt{training.\allowbreak{}input\_\allowbreak{}ranges.\allowbreak{}cone-sec}. \\
\texttt{training.\allowbreak{}torch\_\allowbreak{}version} & string & -- & yes & -- & The torch release the network was trained with. \\
\texttt{training.\allowbreak{}export\_\allowbreak{}parity\_\allowbreak{}max\_\allowbreak{}abs} & number & >= 0 & no & -- & The maximum absolute difference between the trained torch module and the exported float32 graph (run through onnxruntime) on the held-out draw, in standardised output units. \\
\end{longtable}
\endgroup

\subsubsection*{anomaly-model}
% generated by figures/make_tables.py; do not edit by hand
\begingroup\scriptsize
\begin{longtable}{@{}p{0.21\textwidth} p{0.10\textwidth} p{0.19\textwidth} p{0.05\textwidth} p{0.33\textwidth}@{}}
\toprule
Field & Type & Range or values & Req. & Meaning \\
\midrule
\endhead
\bottomrule
\endfoot
\texttt{schema} & const & c\allowbreak{}h\allowbreak{}a\allowbreak{}n\allowbreak{}c\allowbreak{}a\allowbreak{}d\allowbreak{}e\allowbreak{}m\allowbreak{}.\allowbreak{}a\allowbreak{}n\allowbreak{}o\allowbreak{}m\allowbreak{}a\allowbreak{}l\allowbreak{}y\allowbreak{}/\allowbreak{}v\allowbreak{}1 & yes & -- \\
\texttt{engine\_version} & string & pattern \textasciicircum{}\allowbreak{}\textbackslash{}\allowbreak{}d\allowbreak{}+\allowbreak{}\textbackslash{}\allowbreak{}.\allowbreak{}\textbackslash{}\allowbreak{}d\allowbreak{}\{\allowbreak{}2\allowbreak{}\}\allowbreak{}\textbackslash{}\allowbreak{}.\allowbreak{}\textbackslash{}\allowbreak{}d\allowbreak{}\{\allowbreak{}3\allowbreak{}\}\allowbreak{}\$ & yes & -- \\
\texttt{feature\_order} & array & 14 to 14 items & yes & -- \\
\texttt{mean} & array of number & 14 to 14 items & yes & -- \\
\texttt{std} & array of number & 14 to 14 items & yes & -- \\
\texttt{threshold\_p99} & number & > 0 & yes & -- \\
\texttt{train\_recon\_mean} & number & >= 0 & yes & -- \\
\texttt{graph} & object & -- & yes & -- \\
\texttt{graph.path} & const & m\allowbreak{}o\allowbreak{}d\allowbreak{}e\allowbreak{}l\allowbreak{}s\allowbreak{}/\allowbreak{}p\allowbreak{}s\allowbreak{}d\allowbreak{}-\allowbreak{}a\allowbreak{}e\allowbreak{}.\allowbreak{}o\allowbreak{}n\allowbreak{}n\allowbreak{}x & yes & -- \\
\texttt{graph.bytes} & integer & >= 1; <= 5242880 & yes & -- \\
\texttt{graph.sha256} & string & pattern \textasciicircum{}\allowbreak{}[\allowbreak{}0\allowbreak{}-\allowbreak{}9\allowbreak{}a\allowbreak{}-\allowbreak{}f\allowbreak{}]\allowbreak{}\{\allowbreak{}6\allowbreak{}4\allowbreak{}\}\allowbreak{}\$ & yes & -- \\
\texttt{graph.opset} & integer & >= 13 & yes & -- \\
\texttt{graph.input} & const & x & yes & -- \\
\texttt{graph.output} & const & x\allowbreak{}r & yes & -- \\
\texttt{architecture} & string & -- & yes & The layer sizes, for example 14-16-6-16-14, and the input noise used in training. \\
\texttt{training} & object & -- & yes & -- \\
\texttt{training.seed} & integer & -- & yes & -- \\
\texttt{training.n\_train} & integer & >= 1 & yes & -- \\
\texttt{training.noise\_std} & number & >= 0 & yes & -- \\
\texttt{training.input\_ranges} & object & -- & yes & -- \\
\texttt{training.torch\_version} & string & -- & yes & -- \\
\texttt{training.export\_parity\_max\_abs} & number & >= 0 & no & The maximum absolute difference between the trained torch module's reconstruction and the exported graph's (run through onnxruntime) on the ... \\
\end{longtable}
\endgroup

\subsubsection*{model-metrics}
% generated by figures/make_tables.py; do not edit by hand
\begingroup\scriptsize
\begin{longtable}{@{}p{0.20\textwidth} p{0.09\textwidth} p{0.16\textwidth} p{0.04\textwidth} p{0.07\textwidth} p{0.30\textwidth}@{}}
\caption{Fields of the \texttt{model-metrics} contract, version $1.0.1$, generated from its schema file \texttt{model-metrics.schema.json}: the dotted path of each field down to its leaves (\texttt{[]} an element of a list, \texttt{<key>} any key of a map; a structure met again is listed once, where it first appears), its type, its range or allowed values as the schema states them, whether it is required at its level, its unit and its meaning.}\label{tab:schema-model-metrics}\\
\toprule
Field & Type & Range or values & Req. & Unit & Meaning \\
\midrule
\endfirsthead
\multicolumn{6}{@{}l}{\footnotesize Table~\thetable{} (continued)}\\
\toprule
Field & Type & Range or values & Req. & Unit & Meaning \\
\midrule
\endhead
\bottomrule
\endfoot
\texttt{schema} & const & c\allowbreak{}h\allowbreak{}a\allowbreak{}n\allowbreak{}c\allowbreak{}a\allowbreak{}d\allowbreak{}e\allowbreak{}m\allowbreak{}.\allowbreak{}m\allowbreak{}o\allowbreak{}d\allowbreak{}e\allowbreak{}l\allowbreak{}-\allowbreak{}m\allowbreak{}e\allowbreak{}t\allowbreak{}r\allowbreak{}i\allowbreak{}c\allowbreak{}s\allowbreak{}/\allowbreak{}v\allowbreak{}1 & yes & -- & The document schema identifier, chancadem.model-metrics/v1. \\
\texttt{engine\_\allowbreak{}version} & string & pattern \textasciicircum{}\allowbreak{}\textbackslash{}\allowbreak{}d\allowbreak{}+\allowbreak{}\textbackslash{}\allowbreak{}.\allowbreak{}\textbackslash{}\allowbreak{}d\allowbreak{}\{\allowbreak{}2\allowbreak{}\}\allowbreak{}\textbackslash{}\allowbreak{}.\allowbreak{}\textbackslash{}\allowbreak{}d\allowbreak{}\{\allowbreak{}3\allowbreak{}\}\allowbreak{}\$ & yes & -- & The product version that wrote the document (format X.XX.XXX); a version other than VERSION makes the row stale. \\
\texttt{surrogate} & object & -- & yes & -- & The held-out evaluation of the surrogate against the canonical engine. \\
\texttt{surrogate.\allowbreak{}n\_\allowbreak{}train} & integer & >= 1 & yes & -- & The number of training points of the evaluated bundle. \\
\texttt{surrogate.\allowbreak{}n\_\allowbreak{}test} & integer & >= 1 & yes & -- & The number of points of the independent held-out draw. \\
\texttt{surrogate.\allowbreak{}test\_\allowbreak{}seed} & integer & -- & yes & -- & The seed of the held-out Latin-hypercube draw. \\
\texttt{surrogate.\allowbreak{}per\_\allowbreak{}output} & object & -- & yes & -- & The scores of each output, keyed by its name in the bundle's output\_order. \\
\texttt{surrogate.\allowbreak{}per\_\allowbreak{}output.\allowbreak{}<key>.\allowbreak{}r2} & number & <= 1 & yes & -- & Coefficient of determination of the float32 graph against the engine, in original units. \\
\texttt{surrogate.\allowbreak{}per\_\allowbreak{}output.\allowbreak{}<key>.\allowbreak{}mape\_\allowbreak{}pct} & number & >= 0 & yes & \% & MAPE of the float32 graph against the engine, over the points where the engine's value is not zero. \\
\texttt{surrogate.\allowbreak{}per\_\allowbreak{}output.\allowbreak{}<key>.\allowbreak{}mape\_\allowbreak{}pct\_\allowbreak{}int8} & number & >= 0 & yes & \% & The same MAPE for the int8 graph. \\
\texttt{surrogate.\allowbreak{}p80\_\allowbreak{}monotone\_\allowbreak{}vs\_\allowbreak{}css} & boolean & -- & yes & -- & True when the surrogate's P80 does not fall as the closed-side setting opens on every checked curve. \\
\texttt{surrogate.\allowbreak{}parity} & object & -- & yes & -- & The largest differences between the trained module and its exported graphs, in standardised output units. \\
\texttt{surrogate.\allowbreak{}parity.\allowbreak{}torch\_\allowbreak{}vs\_\allowbreak{}float32\_\allowbreak{}max\_\allowbreak{}abs} & number & >= 0 & yes & -- & The trained torch module against the float32 graph, measured by the train stage right after export. \\
\texttt{surrogate.\allowbreak{}parity.\allowbreak{}float32\_\allowbreak{}vs\_\allowbreak{}int8\_\allowbreak{}max\_\allowbreak{}abs} & number & >= 0 & yes & -- & The float32 graph against the int8 graph through onnxruntime on the held-out draw. \\
\texttt{surrogate.\allowbreak{}note} & string & -- & yes & -- & How the surrogate was evaluated, in words. \\
\texttt{anomaly} & object & -- & yes & -- & The held-out evaluation of the autoencoder beside the engine's envelope check. \\
\texttt{anomaly.\allowbreak{}n\_\allowbreak{}heldout} & integer & >= 1 & yes & -- & The number of held-out points. \\
\texttt{anomaly.\allowbreak{}false\_\allowbreak{}alarm\_\allowbreak{}rate} & number & >= 0; <= 1 & yes & fraction & The share of held-out points the autoencoder marks anomalous. \\
\texttt{anomaly.\allowbreak{}n\_\allowbreak{}planted} & integer & >= 1 & yes & -- & The number of planted off-manifold points with finite features. \\
\texttt{anomaly.\allowbreak{}detection\_\allowbreak{}rate} & number & >= 0; <= 1 & yes & fraction & The share of planted points the autoencoder marks anomalous. \\
\texttt{anomaly.\allowbreak{}envelope\_\allowbreak{}false\_\allowbreak{}alarm\_\allowbreak{}rate} & number & >= 0; <= 1 & yes & fraction & The share of held-out points the envelope check flags. \\
\texttt{anomaly.\allowbreak{}envelope\_\allowbreak{}detection\_\allowbreak{}rate} & number & >= 0; <= 1 & yes & fraction & The share of planted points the envelope check flags. \\
\texttt{anomaly.\allowbreak{}planted} & string & -- & yes & -- & How the off-manifold points were planted (which inputs moved outside the training bounds, and by how much). \\
\texttt{anomaly.\allowbreak{}note} & string & -- & yes & -- & The alarm threshold and how the envelope check is read, in words. \\
\end{longtable}
\endgroup

\subsubsection*{readiness}
% generated by figures/make_tables.py; do not edit by hand
\begingroup\scriptsize
\begin{longtable}{@{}p{0.22\textwidth} p{0.11\textwidth} p{0.2\textwidth} p{0.05\textwidth} p{0.34\textwidth}@{}}
\toprule
Field & Type & Range or values & Req. & Meaning \\
\midrule
\endhead
\bottomrule
\endfoot
\texttt{schema} & const & chancadem.readiness/v1 & yes & -- \\
\texttt{product\_version} & string & pattern \textasciicircum{}\textbackslash{}d+\textbackslash{}.\textbackslash{}d\{2\}\textbackslash{}.\textbackslash{}d\{3\}\$ & yes & -- \\
\texttt{catalog\_sha256} & string & pattern \textasciicircum{}[0-9a-f]\{64\}\$ & yes & SHA-256 of docs/contracts/catalog.md (its UTF-8 bytes with CRLF normalised to LF) the document was computed from. \\
\texttt{contracts} & array of object & 1 to n items & yes & -- \\
\texttt{contracts[].id} & string & pattern \textasciicircum{}[a-z0-9][a-z0-9-]*\$ & yes & The catalog row id. \\
\texttt{contracts[].version} & string & pattern \textasciicircum{}\textbackslash{}d+\textbackslash{}.\textbackslash{}d+\textbackslash{}.\textbackslash{}d+\$ & yes & The catalog row's contract version (semver). \\
\texttt{contracts[].status} & enum & real, active & yes & The catalog's design status: real once the real producer serves the row, active while its provider does. \\
\texttt{contracts[].provenance} & enum & real, stub, simulator, fixture, baseline, null & yes & What actually served the row in this build; null when nothing did (state missing). \\
\texttt{contracts[].state} & enum & available, partial, stale, provider, training, missing & yes & -- \\
\texttt{contracts[].reason} & any & -- & yes & Why the state is not available; null only when it is. \\
\texttt{contracts[].missing} & array of string & -- & yes & For partial: what is absent (file paths, rejected record ids). Empty otherwise. \\
\end{longtable}
\endgroup

\subsubsection*{api (OpenAPI routes)}
% generated by figures/make_tables.py; do not edit by hand
\begingroup\scriptsize
\begin{longtable}{@{}p{0.32\textwidth} p{0.12\textwidth} p{0.48\textwidth}@{}}
\caption{Routes of the \texttt{api} contract, version $1.0.1$, generated from the OpenAPI document \texttt{api.openapi.json}: each route with its method, the response codes it declares and its summary.}\label{tab:schema-api}\\
\toprule
Route & Responses & Meaning \\
\midrule
\endfirsthead
\multicolumn{3}{@{}l}{\footnotesize Table~\thetable{} (continued)}\\
\toprule
Route & Responses & Meaning \\
\midrule
\endhead
\bottomrule
\endfoot
\texttt{GET /health} & 200 & Liveness. \\
\texttt{GET /healthz} & 200 & Liveness (alias of /health). \\
\texttt{GET /api/cases} & 200, 503 & The case index (catalog row case-index). \\
\texttt{GET /api/cases/\{case\_id\}/manifest} & 200, 404, 503 & One case manifest (catalog row case-manifest). \\
\texttt{GET /api/artifacts/\{rel\_path\}} & 200, 404, 503 & An artifact a manifest names (catalog row trace), by its path under data/derived. \\
\texttt{GET /api/readiness} & 200, 503 & What is real in this build (catalog row readiness). \\
\end{longtable}
\endgroup

